\section{Scope, conventions, and the status of the results}
\label{sec:scope}

The point of departure is the 228-page collective manuscript in
\path{Analysis/Polylogarithms/docs/manuscript} of the ProveIt repository
\cite{ProveIt}. The source was pinned before starting this continuation:
\begin{center}
\small\texttt{9bc738d3be22b8586a24693f19fb2e2a50ecd1bf}.
\end{center}
This matters because the manuscript already contains several earlier
continuations. Its Gaussian parity reductions, depth-four reciprocal
identities, modular row convergence bounds, and global Stieltjes zero
theorem are established material. They are used as context or reproved
where needed, rather than presented as new discoveries.

The present work follows three distinct threads. The first concerns the
exact integer structure of a specified relation matrix. The second turns
empirical CM products into consequences of the arithmetic of modular
functions and a classical period formula. The third determines positive
real zeros by combining global analytic bounds with finite certificates.
These threads share a methodological point: a proposed identity or zero
count becomes substantially more useful when its precise proof mechanism
and its limits are stated.

\subsection{What is proved, and what remains open}

The following map separates the main claims and their evidence.
\begin{center}
\small
\begin{tabularx}{\textwidth}{@{}>{\raggedright\arraybackslash}p{0.24\textwidth}>{\raggedright\arraybackslash}X>{\raggedright\arraybackslash}p{0.23\textwidth}@{}}
\toprule
Claim & Result and proof mechanism & Location\\
\midrule
Odd-weight formal ranks & The rank conjecture holds for every odd weight;
integer polynomial substitutions also determine saturation of the row
lattice. & Theorem~\ref{newrank:thm:smith}\\[3pt]
Even-weight formal system & Full column rank; all nonunit invariant
factors equal $2$. & Theorem~\ref{newrank:thm:smith}\\[3pt]
Mixed-color values & A finite product certificate evaluates
$\Im\Li_{w-1,1}(i,-1)$ at every even weight. &
Theorem~\ref{newmixed:thm:family}\\[3pt]
CM class products & An all-even-weight resultant formula; five baseline
weight-four products and three displayed weight-six extensions. &
Theorem~\ref{cmnorm:thm:allweights}\\[3pt]
CM genus ratios & The recorded ratios and weight-six extensions are
proved; one number field of degree at most $12$ contains every even-weight
ratio for a fixed character. & Theorems~\ref{cmgenus:thm:four},
\ref{cmgenus:thm:degree}\\[3pt]
Fifth-index real zeros & Exactly three at derivative orders one and two,
and five thereafter, all simple. & Theorem~\ref{n5:thm:transition}\\[3pt]
Lerch deformation & Complete small-parameter classification, alternating
branch directions, and a proved double-zero bifurcation. &
Theorems~\ref{thm:small-rho}, \ref{thm:rho-direction}, \ref{thm:fold}\\
\bottomrule
\end{tabularx}
\end{center}

The Gaussian matrix theorem is about \emph{formal relations}. It does not
compute the dimension of a space of actual transcendental numbers. The CM
identities use classical period theory; the contribution is a fully
specified formula in the manuscript's normalization and its exact
applications. The fifth-index zero count uses a computer-assisted finite
lemma with explicit rational arithmetic, not an unproved numerical root
scan. The small-parameter and bifurcation results have analytic proofs;
decimal bifurcation coordinates are supplementary diagnostics.

We do not settle the manuscript's $S_4$ mixed-relation conjecture, general
rational-argument weight-four ladders, or its highest-weight algebraic
ladder claims. In particular, an obstruction to one finite relation system
is not an obstruction to all functional equations. Priority claims here
are relative to the pinned manuscript. The literature survey identifies
the classical ingredients and relevant existing methods, but does not
justify asserting that every displayed consequence is absent from the
entire historical literature.

\subsection{Polylogarithms, convergence, and color conventions}

We use the strict nested-sum convention
\begin{equation}\label{eq:depth-two-convention}
 \Li_{a,b}(x,y)=\sum_{m>n\ge1}\frac{x^m y^n}{m^a n^b},
 \qquad \Li_s(z)=\sum_{m\ge1}\frac{z^m}{m^s}.
\end{equation}
This fixes the distinction between colors $(x,y)$ and the adjacent-letter
coordinates in an iterated integral. At roots of unity the values are
ordinary convergent boundary values where the series converges; the one
divergent imaginary coordinate that occurs in the formal system is
explicitly omitted. Its regularized cancellation rows are exactly those
specified by the baseline. No additional regularization convention is
silently introduced.

For positive integer $s$, write
\[
 \beta(s)=\sum_{m\ge0}\frac{(-1)^m}{(2m+1)^s}
          =\Im\Li_s(i),\qquad
 \eta_{\mathrm{alt}}(s)=\sum_{m\ge1}\frac{(-1)^{m-1}}{m^s}.
\]
Thus $\eta_{\mathrm{alt}}(1)=\log2$ and
$\eta_{\mathrm{alt}}(s)=(1-2^{1-s})\zeta(s)$ for $s>1$.
The subscript distinguishes this alternating Dirichlet series from the
Dedekind eta function $\eta(\tau)$ in the CM sections.

The colored shuffle and stuffle framework is standard. Zhao's work
\cite{Zhao} distinguishes the various classes of standard relations,
while Yuan--Zhao \cite{YuanZhao} provide closely relevant formal
double-space and polynomial constructions. We use this context to locate
the exact matrix theorem, without identifying its particular quotient
with a larger motivic or numerical period space.

\subsection{Analytic conventions}

The generalized Stieltjes constants are defined by
\begin{equation}\label{eq:laurent-convention}
 \zeta(s,a)=\frac1{s-1}+\sum_{n\ge0}
       \frac{(-1)^n}{n!}\gamma_n(a)(s-1)^n,
 \qquad a>0.
\end{equation}
A superscript $(k)$ on $\gamma_n$ means an $a$ derivative. A derivative
$\partial_s$ of a zeta or Lerch function always means a derivative in its
complex order. We keep these two operations separate. The elementary
Lerch endpoint is $\rho=0$; the classical Hurwitz endpoint is $\rho=1$.
The coefficients themselves need not have a convergent series at the
classical endpoint, but their parameter derivatives of positive order do.
Definitions and the standard Mellin representations are recorded in the
DLMF \cite{DLMFHurwitz,DLMFLerch}.

All root counts are on the open positive real axis. A count with
\emph{multiplicity} is explicitly identified as such. When a theorem
asserts simplicity, the multiplicity count and the number of distinct
zeros agree. This distinction is essential at the elementary endpoint
and at a bifurcation.
